% !PRINTED 502
\weberpage{502}

By § 38 (15) we therefore have

\weq{10}{M = (-1)^{A + C + \frac{B+1}{2}} \sqrt{D}^3\, \gamma_3(\omega),}
and here there occurs only one of the values $v = u$; accordingly $M$,
and consequently also $\sqrt{D}\,\gamma_3(\omega)$, is rationally expressible through $j(\omega)$.

In order to dispose of the exceptional case as well, we set
\[
n = -D = 3\,m^2.
\]
Then (9) becomes:
\[
y^2 + 3\,m^2(x^2 - 4) = 0,
\]
and this has the three solutions:
\[
\begin{aligned}
x &= 2, & y &= 0, \\
x &= 1, & y &= \mp 3\,m.
\end{aligned}
\]
For $x = 2$ the equations (8) again give the formula (10):
\[
M_1 = (-1)^{A + C + \frac{B+1}{2}} \sqrt{D}^3\, \gamma_3(\omega),
\]
and the two other cases give
\[
\begin{aligned}
\partial &= 1, & a &= n, \\
\alpha &= \frac{B \pm 3\,m}{2}, & \beta &= A, \\
\gamma &\equiv C, & \alpha + \delta &\equiv B \pmod{4}, \\
\alpha + \beta\omega &= \frac{\sqrt{D} \pm 3\,m}{2}. &&
\end{aligned}
\]
From this one obtains by § 38 (15) for $M_2$, $M_3$ the values
\[
M_2 = (-1)^{\frac{B-1}{2}} \left(\frac{3\,m + \sqrt{D}}{2}\right)^3 \gamma_3(\omega),
\]
\[
M_3 = (-1)^{\frac{B-1}{2}} \left(\frac{-3\,m + \sqrt{D}}{2}\right)^3 \gamma_3(\omega).
\]
Now it is true that the quantities $M_1$, $M_2$, $M_3$ are not rationally
expressible through $j(\omega)$ individually, but, as we shall show at once, their symmetric functions are, for example their sum, and
from this it follows for this case also that $\sqrt{D}\,\gamma_3(\omega)$ is a class invariant.

From the assumption we have made, that $A$ is relatively
prime to $2\,D$, we can free ourselves afterwards, since $\gamma_3(\omega)$
changes at most its sign under all linear transformers.
We then have the theorem:

1. \gesperrt{If $D \equiv 1 \pmod{4}$, then $\sqrt{D}\,\gamma_3(\omega)$ is a class invariant of the discriminant $D$.}
